%!TEX program = pdflatex
\documentclass[UTF8]{elegantpaper}
\usepackage{ctex}
\usepackage{geometry}
\usepackage{ulem}
\usepackage{calc}
\usepackage{tikz}
\geometry{left=3.18cm,right=3.18cm,top=2.54cm,bottom=2.54cm}

\newlength{\textlen}
\newlength{\lablen}
\setlength{\lablen}{0pt}

\newcommand\college{学院}
\newcommand\majorID{专业}
\newcommand\IDnumber{学号}
\newcommand\thesisauthor{姓名}
\newcommand\adviser{指导教师}
\newcommand\coadviser{协助指导教师}
\newcommand\applydate{日期}

\begin{document}
\thispagestyle{empty}
\begin{figure}[t]
	\parbox[b]{2cm}{
		\includegraphics[scale=0.22]{HIT.jpg} % 校徽图片
	}
	\parbox[b]{13.5cm}{
		\begin{center}
			\includegraphics[scale=0.8]{hitname.jpg} 
			% 学校名图片，两张图片必须放在与.tex文件同一目录下，名字可以自己相应更改，但必须保持一致
			
			\large \textbf{Harbin Institute of Technology} 
		\end{center}
	}
\end{figure}

\begin{center}
	\quad \\
	\quad \\
	\quad \\
	\heiti \fontsize{30}{12pt} 大标题 % 大标题 黑体，字号30，行距12pt
	\vskip 2.5cm
	\heiti \zihao{2} 小标题 %小标题可插入 \\ 换行，黑体，字号2
\end{center}
\vskip 6cm

\begin{center}
	\foreach \x in {\college, \majorID, \thesisauthor, \adviser, \coadviser, \applydate} {
		\settowidth{\textlen}{\x}
		\ifdim \textlen > \lablen
		\setlength{\lablen}{\textlen}
		\else
		\relax
		\fi
	}
	\setlength{\lablen}{\widthof{\coadviser}}
	\begin{tabular} {ll}
		% 如果不需要某一行，直接将整行删去或注释即可
		
		% 学院行
		\makebox[\the\lablen][s]{\college} ：     & \underline{\makebox[4cm][c]{XXXXXXX学院}} 		\\ 
		
		% 专业行
		\makebox[\the\lablen][s]{\majorID} ：     & \underline{\makebox[4cm][c]{XX专业}} 		\\ 
		
		% 学号行
		\makebox[\the\lablen][s]{\IDnumber} ：    & \underline{\makebox[4cm][c]{1XXXXXXXXX}} 		\\ 
		
		% 姓名行
		\makebox[\the\lablen][s]{\thesisauthor} ：& \underline{\makebox[4cm][c]{姓名}} 		\\ 
		
		% 指导教师行
		\makebox[\the\lablen][s]{\adviser} ：     & \underline{\makebox[4cm][c]{指导教师}} 		\\ 
		
		% 协助指导教师行
		\makebox[\the\lablen][s]{\coadviser} ：   & \underline{\makebox[4cm][c]{协助指导教师}} 		\\ 
		
		% 日期行
		\makebox[\the\lablen][s]{\applydate} ：   & \underline{\makebox[4cm][c]{日期}} 		\\ 	
	\end{tabular}
	\vspace{4ex}
	
\end{center}

% \newpage 使用时取消此行注释即可添加新界面
\end{document}